\section{Results}
\label{sec:results}

\subsection{System Performance and Reliability}

The SAGE framework demonstrates exceptional system-level performance across 600 evaluations. Table~\ref{tab:system_reliability} summarizes key reliability metrics.

\begin{table}[t]
\centering
\caption{System reliability and quality assurance metrics}
\label{tab:system_reliability}
\small
\resizebox{\columnwidth}{!}{%
\begin{tabular}{lcc}
\toprule
\textbf{Metric} & \textbf{Value} & \textbf{Interpretation} \\
\midrule
\multicolumn{3}{l}{\textit{Evaluation Success}} \\
Total evaluations & 600 & Perfect dataset balance \\
Success rate & 100.0\% & All evaluations completed \\
Parsing success & 100.0\% & All outputs valid JSON \\
\midrule
\multicolumn{3}{l}{\textit{Convergence \& Consistency}} \\
Convergence rate & 98.8\% & Score stability (R4$\rightarrow$R5) \\
Peer agreement & $>$94\% & Inter-rater reliability \\
Mean IRR difference & $<$0.3 & Excellent agreement \\
\midrule
\multicolumn{3}{l}{\textit{Mode Robustness}} \\
Mode invariance & Max $|\Delta|$=0.05 & Content $\approx$ Title \\
\bottomrule
\end{tabular}
}
\vspace{0.1cm}

\small
\textit{Notes:} Convergence = absolute score change $<$0.3 from Round 4 to Round 5. Peer agreement = primary-validator difference $<$0.5. IRR = inter-rater reliability. Mode invariance = maximum absolute difference between content-limit and title-limit across all layers.
\end{table}

All 600 evaluations completed successfully with valid JSON outputs, achieving a 100\% success rate. The five-round iterative reflection mechanism achieves 98.8\% convergence rate, with dimension scores stabilizing between Round 4 and Round 5. Figure~\ref{fig:convergence} illustrates the score convergence trajectories across all five rounds for each layer and genre category, showing rapid stabilization by Rounds 3--4. Independent validation through peer review yields greater than 94\% agreement rate across all dimensions and layers, with mean inter-rater difference below 0.3 points.

\begin{figure*}[t]
\centering
\includegraphics[width=\textwidth]{sections/figures/convergence_trajectory.png}
\caption{Score convergence trajectories across five iterative rounds for each layer (L4--L6) and genre category. Scores stabilize rapidly by Round 3--4, with canonical and pulp fiction consistently scoring higher than LLM-generated stories across all layers and rounds.}
\label{fig:convergence}
\end{figure*}

Table~\ref{tab:mode_invariance} presents mode comparison across all three layers, with maximum absolute difference of 0.05 between content-limit and title-limit modes. Paired samples t-tests confirm no significant mode effects (all $p>0.05$).

\begin{table}[t]
\centering
\caption{Mode invariance: Content-limit vs Title-limit comparison}
\label{tab:mode_invariance}
\small
\resizebox{\columnwidth}{!}{%
\begin{tabular}{lccc}
\toprule
\textbf{Layer} & \textbf{Content-Limit} & \textbf{Title-Limit} & \textbf{Difference} \\
\midrule
Layer 4 (Cultural)     & 3.63 & 3.65 & +0.02 \\
Layer 5 (Emotional)    & 3.98 & 3.93 & $-$0.05 \\
Layer 6 (Existential)  & 3.60 & 3.59 & $-$0.01 \\
\midrule
\textbf{Overall}       & \textbf{3.74} & \textbf{3.72} & \textbf{$-$0.01} \\
\bottomrule
\end{tabular}
}
\vspace{0.1cm}

\small
\textit{Notes:} Mean scores across 100 stories per mode. All differences $<$0.05 indicate mode-invariant performance. Content-limit: evaluator receives only story text (no metadata). Title-limit: evaluator receives only title + author (no text).
\end{table}

\subsection{Genre Discrimination}

Figure~\ref{fig:genre_comparison} presents genre comparison across all three analytical layers, revealing a consistent and statistically significant quality hierarchy. Detailed numerical results with standard deviations and pairwise differences are provided in Table~\ref{tab:genre_comparison_appendix} (Appendix).

\begin{figure}[t]
\centering
\includegraphics[width=\columnwidth]{sections/figures/genre_comparison.png}
\caption{Genre comparison across three analytical layers. Error bars indicate standard deviations. All pairwise differences between canonical and LLM-generated stories are significant at $p<0.001$.}
\label{fig:genre_comparison}
\end{figure}

Across all three layers, canonical literature significantly outperforms both pulp fiction and LLM-generated stories (all pairwise $p<0.001$), yielding a consistent genre ranking: Canonical (4.02) $>$ Pulp (3.85) $>$ LLM (2.83). The Canonical vs.\ LLM difference of +1.19 points (29.6\% relative gap) with very large effect sizes (Cohen's $d>1.5$ across all layers) confirms that the framework identifies fundamental quality differences between human-authored canonical literature and AI-generated narratives. Figure~\ref{fig:effect_sizes} visualizes these effect sizes across layers. The Canonical vs.\ Pulp difference remains modest (+0.17 points, 4.4\%).

\begin{figure}[t]
\centering
\includegraphics[width=\columnwidth]{sections/figures/effect_sizes.png}
\caption{Effect size (Cohen's $d$) for Canonical vs.\ LLM-generated stories across three layers. All values exceed the ``large effect'' threshold ($d=0.8$), indicated by the dashed red line.}
\label{fig:effect_sizes}
\end{figure}

\subsection{Layer-Specific Results}

\subsubsection{Layer 4: Cultural Representation}

Layer 4 achieves a mean score of 3.64 across 200 evaluations (100 stories $\times$ 2 modes). Canonical works (3.96 $\pm$ 0.41) significantly outperform pulp fiction (3.83 $\pm$ 0.36, $p<0.01$) and LLM-generated stories (2.55 $\pm$ 0.62, $p<0.001$), yielding the largest Canonical vs.\ LLM difference among the three layers (+1.41 points, 35.6\% relative gap, $d$=2.68). All four cultural dimensions (Intersectional Power Dynamics, Cultural Voice \& Perspective, Cultural Specificity, and Cultural Pattern Complexity) exhibit consistent genre hierarchy with uniformly large effect sizes.

\subsubsection{Layer 5: Emotional-Psychological Representation}

Layer 5 achieves the highest mean score (3.96) across all three layers. Canonical literature (4.15 $\pm$ 0.31) significantly outperforms LLM-generated stories (3.36 $\pm$ 0.59, $p<0.001$, $d$=1.68), but the 19.1\% relative gap is the smallest among the three layers. Pulp fiction (4.04 $\pm$ 0.24) approaches canonical performance with only 0.11-point difference (2.7\%), the smallest Canonical vs.\ Pulp gap across all layers.

\subsubsection{Layer 6: Existential-Philosophical Representation}

Layer 6 demonstrates strong genre discrimination comparable to Layer 4. Canonical works (3.95 $\pm$ 0.57) substantially exceed LLM-generated stories (2.59 $\pm$ 0.57) by +1.36 points (34.5\% relative gap, $p<0.001$, $d$=2.40). The Canonical vs.\ Pulp difference (+0.27, 6.8\%) is the largest among the three layers, with pulp fiction (3.68 $\pm$ 0.56) achieving moderate existential engagement.

\subsection{Cross-Layer Correlation}

To assess whether the three analytical layers capture genuinely distinct quality dimensions (RQ4), we compute pairwise correlation coefficients across all 200 story-mode observations. Table~\ref{tab:cross_layer_correlation} presents the results.

\begin{table}[t]
\centering
\caption{Cross-layer correlation analysis (Pearson $r$)}
\label{tab:cross_layer_correlation}
\small
\begin{tabular}{lccc}
\toprule
\textbf{Layer Pair} & \textbf{Pearson $r$} & \textbf{Spearman $\rho$} & \textbf{$p$-value} \\
\midrule
L4 $\times$ L5 (Cultural $\leftrightarrow$ Emotional)     & 0.649 & 0.514 & $<$0.001 \\
L4 $\times$ L6 (Cultural $\leftrightarrow$ Existential)   & 0.683 & 0.598 & $<$0.001 \\
L5 $\times$ L6 (Emotional $\leftrightarrow$ Existential)  & 0.680 & 0.621 & $<$0.001 \\
\bottomrule
\end{tabular}
\vspace{0.1cm}

\small
\textit{Notes:} $N$=200 (100 stories $\times$ 2 modes). Spearman rank correlations provided as robustness check. All correlations significant at $p<0.001$.
\end{table}

All three layer pairs exhibit moderate positive correlations (Pearson $r$ = 0.649--0.683), statistically significant at $p<0.001$, yet falling below the $r>0.7$ threshold that would suggest dimensional redundancy. Spearman rank correlations ($\rho$ = 0.514--0.621) corroborate this pattern. The three layers exhibit distinct overall means: Layer 5 (Emotional, 3.96) $>$ Layer 4 (Cultural, 3.64) $>$ Layer 6 (Existential, 3.59), indicating that emotional depth is the most widespread quality dimension while cultural critique and philosophical engagement vary more substantially by genre and authorship type.
